%=======================================================================
%  IE 5311-001 (D01) Optional Course Project
%  CVRP: MIP, Dantzig--Wolfe decomposition and column generation
%  Instance E-n22-k4 (CVRPLIB). At most 10 pages.
%  Packages: geometry, amsmath, amssymb, booktabs, array.
%=======================================================================
\documentclass[11pt]{article}

\usepackage[margin=1in]{geometry}
\usepackage{amsmath}
\usepackage{amssymb}
\usepackage{booktabs}
\usepackage{array}

\DeclareMathOperator*{\minimize}{min}
\DeclareMathOperator{\conv}{conv}
\newcommand{\T}{^{\!\top}}
\newcommand{\R}{\mathbb{R}}

\setlength{\parskip}{0.35em}
\setlength{\parindent}{0pt}
\setlength{\abovedisplayskip}{5pt}
\setlength{\belowdisplayskip}{5pt}

\title{\vspace{-1.2cm}Capacitated Vehicle Routing by Dantzig--Wolfe
       Decomposition and Column Generation\\
       \large IE 5311-001 (D01) Optional Course Project}
\author{Lidivine Kengne}
\date{Fall 2026}

\begin{document}
\maketitle
\vspace{-0.9cm}

\section{Data and Preliminary Analysis}

\paragraph{(a) Notation.}
Let $V = \{1,\dots,n\}$ be the node set with $n = 22$, the depot being
node $o = 1$ and $C = V \setminus \{o\}$ the $21$ customers. Let
$\mathcal{K} = \{1,\dots,K\}$ index the identical vehicles, $K = 4$, and
$A = \{(i,j) : i,j \in V,\; i \neq j\}$ the arc set. Each customer
$i \in C$ has demand $d_i > 0$ (with $d_o = 0$), each vehicle has
capacity $Q = 6000$, and each arc carries cost
\[
  c_{ij} \;=\;
  \Bigl\lfloor \sqrt{(x_i - x_j)^2 + (y_i - y_j)^2} + 0.5 \Bigr\rfloor ,
  \qquad (i,j) \in A,
\]
the CVRPLIB \texttt{EUC\_2D} convention: Euclidean distance between the
node coordinates rounded to the nearest integer. The cost matrix is
symmetric and satisfies the triangle inequality.

\paragraph{(b) Total demand and a fleet lower bound.}
Summing the demands gives
\[
  D \;=\; \sum_{i \in C} d_i \;=\; 22{,}500 .
\]
Any feasible solution assigns every customer to exactly one vehicle and
no vehicle may carry more than $Q$, so the fleet must satisfy
$K \cdot Q \ge D$, giving the bound
\[
  K \;\ge\; \Bigl\lceil \frac{D}{Q} \Bigr\rceil
  \;=\; \Bigl\lceil \frac{22{,}500}{6{,}000} \Bigr\rceil
  \;=\; \lceil 3.75 \rceil \;=\; 4 .
\]
The instance name specifies $K = 4$, so the bound is tight and the
fleet has only $KQ - D = 24{,}000 - 22{,}500 = 1{,}500$ units of total
slack, about $6.25\%$ of capacity.

\emph{Does meeting the bound guarantee feasibility?} No. The bound is
necessary, not sufficient. It treats demand as perfectly divisible,
whereas a customer's demand cannot be split, so the real question is
whether the $21$ indivisible demands can be packed into $4$ bins of
size $Q$ --- a bin-packing feasibility problem, itself NP-complete.
Aggregate capacity can be ample while no packing exists: three
customers of demand $2{,}100$ each and a fleet of $K=1$, $Q=6{,}000$
satisfies $KQ \ge D$ with $D = 6{,}300 > Q$ failing already, and more
pointedly $d = (4{,}000,\,4{,}000)$ with $K = 2$, $Q = 6{,}000$ gives
$KQ = 12{,}000 \ge 8{,}000 = D$ yet forces one vehicle per customer,
which is feasible, while $d = (4{,}000,\,4{,}000,\,4{,}000)$ with
$K = 2$ has $KQ = 12{,}000 = D$ and no packing at all. For this
instance a feasible packing does exist, but that is established by
exhibiting a solution, not by the counting bound.

\paragraph{(c) Data check.}
With CVRPLIB numbering, node $1$ is the depot at $(145,215)$, node $2$
is at $(151,264)$ and node $3$ at $(159,261)$:

\begin{center}
\begin{tabular}{@{}lllrr@{}}
\toprule
arc & from & to & exact distance & $c_{ij}$\\
\midrule
$(1,2)$ & $(145,215)$ & $(151,264)$ & 49.365980 & 49\\
$(1,3)$ & $(145,215)$ & $(159,261)$ & 48.083261 & 48\\
$(2,3)$ & $(151,264)$ & $(159,261)$ &  8.544004 &  9\\
\bottomrule
\end{tabular}
\end{center}

The demands run from $100$ to $2{,}500$; since $\max_i d_i = 2500 \le Q$
every customer is individually servable, so no instance-level
infeasibility arises from a single oversized demand.

\section{MIP Formulation and Dantzig--Wolfe Decomposition}

\subsection{(a) Original vehicle-indexed MIP}

\paragraph{Decision variables.} For $(i,j) \in A$ and $k \in \mathcal{K}$,
\[
  x_{ijk} =
  \begin{cases}
    1, & \text{vehicle } k \text{ traverses arc } (i,j),\\
    0, & \text{otherwise,}
  \end{cases}
  \qquad\text{and write}\qquad
  y_{ik} \;=\; \sum_{j \in V \setminus \{i\}} x_{jik}
\]
for the derived indicator that vehicle $k$ visits node $i$. The
formulation is
\begin{align}
  \minimize \quad & \sum_{k \in \mathcal{K}} \sum_{(i,j) \in A}
      c_{ij}\, x_{ijk} \label{eq:obj}\\
  \text{s.t.}\quad
  & \sum_{k \in \mathcal{K}} y_{ik} = 1,
    && i \in C, \label{eq:visit}\\
  & \sum_{j \neq h} x_{jhk} \;=\; \sum_{j \neq h} x_{hjk},
    && h \in V,\; k \in \mathcal{K}, \label{eq:flow}\\
  & \sum_{j \in C} x_{ojk} = 1,
    && k \in \mathcal{K}, \label{eq:depart}\\
  & \sum_{i \in C} d_i\, y_{ik} \;\le\; Q,
    && k \in \mathcal{K}, \label{eq:cap}\\
  & \sum_{i \in S} \sum_{j \in S \setminus \{i\}} x_{ijk}
      \;\le\; |S| - 1,
    && S \subseteq C,\; |S| \ge 2,\; k \in \mathcal{K},
    \label{eq:sec}\\
  & x_{ijk} \in \{0,1\},
    && (i,j) \in A,\; k \in \mathcal{K}. \label{eq:bin}
\end{align}

\paragraph{How the requirements are enforced.}
(i) Constraint \eqref{eq:visit} states that across all vehicles node
$i$ is entered exactly once, so every customer is served exactly once
and demands are never split. (ii) Constraint \eqref{eq:depart} forces
each vehicle to leave the depot exactly once, towards a customer, so
its route is nonempty; \eqref{eq:flow} balances in-degree against
out-degree at every node, which forces the vehicle to return to the
depot and makes its arc set a union of closed walks. (iii) Constraint
\eqref{eq:cap} charges each vehicle the demand of the nodes it visits
and caps it at $Q$. (iv) Constraint \eqref{eq:sec} is the
Dantzig--Fulkerson--Johnson family: any customer subset $S$ may carry
at most $|S|-1$ internal arcs of a single vehicle, which forbids a
closed walk contained entirely in $C$. Together with \eqref{eq:flow}
and \eqref{eq:depart} this leaves exactly one cycle per vehicle, and it
passes through the depot. The formulation has $|A| \cdot K = 1{,}848$
binaries; \eqref{eq:sec} is exponential in $|C|$ and is therefore
imposed lazily, as described in Section~3.

\paragraph{Classification (slides 8--10).} Constrained; discrete
variables, huge but finite; linear objective and constraints;
deterministic; single stage; single player.

\subsection{(b) Dantzig--Wolfe reformulation}

\paragraph{Which constraints couple.} Only \eqref{eq:visit} involves
more than one vehicle: it is the single family tying the $K$ vehicles
together. Constraints \eqref{eq:flow}--\eqref{eq:bin} are separable by
$k$: for each vehicle they describe that vehicle's own feasible routes
and say nothing about the others. This is exactly the block-angular
structure Dantzig--Wolfe exploits.

\paragraph{The subproblem.} For vehicle $k$ define
\[
  X_k \;=\; \bigl\{\, x^k \in \{0,1\}^{A} \;:\;
    \eqref{eq:flow},\,\eqref{eq:depart},\,\eqref{eq:cap},\,\eqref{eq:sec}
    \text{ hold for } k \,\bigr\} .
\]
By the argument above, $X_k$ is precisely the set of incidence vectors
of \emph{feasible single-vehicle routes}: elementary cycles through the
depot whose total demand does not exceed $Q$. It is a finite set of
$0$--$1$ points, so $\conv(X_k)$ is a polytope whose extreme points are
the elements of $X_k$ themselves. Enumerate them as
$x^{k,1},\dots,x^{k,R_k}$.

\paragraph{Substitution.} Every $x^k \in \conv(X_k)$ is a convex
combination $x^k = \sum_{r} \lambda_{kr} x^{k,r}$ with
$\sum_r \lambda_{kr} = 1$, $\lambda_{kr} \ge 0$. Writing
$c_{kr} = \sum_{(i,j)} c_{ij} x^{k,r}_{ij}$ for the cost of route $r$
and $a_{i}^{kr} = \sum_{j \neq i} x^{k,r}_{ji} \in \{0,1\}$ for whether
it visits customer $i$, substitution into \eqref{eq:obj} and
\eqref{eq:visit} gives the master
\begin{align}
  \minimize \quad & \sum_{k \in \mathcal{K}} \sum_{r=1}^{R_k}
      c_{kr}\, \lambda_{kr}\\
  \text{s.t.}\quad
  & \sum_{k \in \mathcal{K}} \sum_{r=1}^{R_k}
      a_i^{kr}\, \lambda_{kr} = 1, && i \in C,\\
  & \sum_{r=1}^{R_k} \lambda_{kr} = 1, && k \in \mathcal{K},\\
  & \lambda_{kr} \ge 0 .
\end{align}

\paragraph{Equivalence under integrality.} Impose
$\lambda_{kr} \in \{0,1\}$. The convexity row then selects exactly one
index $r$ per vehicle, so the recovered $x^k = x^{k,r}$ lies in $X_k$,
and the coupling row says the chosen routes partition $C$. Conversely
any feasible point of \eqref{eq:obj}--\eqref{eq:bin} has each $x^k \in
X_k$, hence equals some extreme point, giving a $0$--$1$ $\lambda$ of
the same cost. The two integer programs therefore have the same
feasible solutions and the same objective values: they are equivalent.
The difference is in their \emph{relaxations}, which is the subject of
Section~5.

\subsection{(c) Identical vehicles: a route-based set-partitioning model}

All $K$ vehicles share the same capacity $Q$ and the same cost matrix,
so $X_k = X$ for every $k$: the feasible route set and the route costs
do not depend on $k$. Let $\mathcal{R}$ enumerate the feasible routes
once, with cost $c_r$ and coverage $a_{ir} \in \{0,1\}$. Aggregating
the vehicle-indexed weights,
\[
  \lambda_r \;=\; \sum_{k \in \mathcal{K}} \lambda_{kr},
\]
and summing the $K$ convexity rows gives
$\sum_{r \in \mathcal{R}} \lambda_r = K$. The master collapses to
\begin{align}
  \minimize \quad & \sum_{r \in \mathcal{R}} c_r \lambda_r
    \label{eq:spobj}\\
  \text{s.t.}\quad
  & \sum_{r \in \mathcal{R}} a_{ir} \lambda_r = 1, && i \in C,
    \label{eq:spcov}\\
  & \sum_{r \in \mathcal{R}} \lambda_r = K, \label{eq:spfleet}\\
  & \lambda_r \in \{0,1\} \quad\text{(or } \lambda_r \ge 0
    \text{ in the LP relaxation)}. \label{eq:spint}
\end{align}
Constraint \eqref{eq:spcov} is the customer-covering family: it is the
aggregated coupling constraint \eqref{eq:visit}, and as an equality it
enforces exactly-once service, making \eqref{eq:spobj}--\eqref{eq:spint}
a set-\emph{partitioning} model. Constraint \eqref{eq:spfleet}
represents the fixed fleet: it is the sum of the $K$ separate convexity
rows, which is all that survives once the vehicles are interchangeable.
Binary rather than general integer $\lambda_r$ loses nothing, since
$\lambda_r \ge 2$ would serve every customer of $r$ at least twice and
violate \eqref{eq:spcov}.

The practical gain is the removal of the vehicle index. The original
formulation is symmetric under all $K! = 24$ relabellings of identical
vehicles, so a branch-and-bound search re-examines each solution up to
$24$ times; the route-based model has no such symmetry.

\subsection{(d) Column generation}

$|\mathcal{R}|$ grows exponentially with $|C|$, so the LP relaxation of
\eqref{eq:spobj}--\eqref{eq:spint} cannot be written down in general.
Column generation solves it while handling only a small subset.

\paragraph{Restricted master.} Choose
$\mathcal{R}' \subseteq \mathcal{R}$ and solve
\eqref{eq:spobj}--\eqref{eq:spfleet} over $\mathcal{R}'$ with
$\lambda \ge 0$. Let $\pi_i$ be the dual of the covering row for
customer $i$ and $\sigma$ the dual of the fleet row.

\paragraph{Reduced cost and the pricing subproblem.} A column
$r \in \mathcal{R}$ has reduced cost
\[
  \bar{c}_r \;=\; c_r \;-\; \sum_{i \in C} a_{ir}\pi_i \;-\; \sigma ,
\]
so the most attractive column solves
\begin{equation}\label{eq:pricing}
  \minimize_{r \in \mathcal{R}}\; \bar{c}_r
  \;=\; -\sigma \;+\;
  \minimize_{r \in \mathcal{R}}
  \Bigl\{ \sum_{(i,j) \in r} c_{ij} - \sum_{i \in r} \pi_i \Bigr\} .
\end{equation}
Charging each dual at the head of the arc, $\tilde{c}_{ij} =
c_{ij} - \pi_j$ with $\pi_o = 0$, makes \eqref{eq:pricing} a shortest
closed walk from the depot to itself under modified costs, subject to
visiting each customer at most once and to the load limit: an
\emph{elementary shortest path problem with resource constraints}. It
is the single-vehicle problem $X$ again, with a different objective.
Explicitly, with $x_{ij}$ and $y_h = \sum_i x_{ih}$,
\[
  \minimize \;\; \sum_{(i,j) \in A} c_{ij} x_{ij}
      - \sum_{h \in C} \pi_h y_h
  \quad\text{s.t.}\quad
  \sum_{j \in C} x_{oj} = 1,\;\;
  \sum_{i \in C} x_{io} = 1,\;\;
  y_h = \!\!\sum_{j \neq h}\!\! x_{hj} \le 1,\;\;
  \sum_{h \in C} d_h y_h \le Q,
\]
together with \eqref{eq:sec} for the single vehicle and
$x \in \{0,1\}^A$.

\paragraph{Procedure.}
\emph{Initialisation.} $\mathcal{R}'$ must make the restricted master
feasible, which needs $K$ routes partitioning $C$. We obtain them from
a first-fit-decreasing packing of the demands into $K$ bins followed by
an exact traveling-salesman ordering within each bin, so the master is
feasible at the first iteration and no artificial variables or Phase~I
are required. \emph{Iteration.} Solve the restricted master, read
$(\pi,\sigma)$, solve \eqref{eq:pricing}; if $\bar{c}_r < 0$ append the
route and repeat. \emph{Termination.} Stop when the pricing optimum
satisfies $\bar{c}_r \ge 0$. Because the pricing problem is solved
\emph{exactly} over all of $\mathcal{R}$, this certifies that no column
of the full master can improve the current basis, so the restricted
master optimum is the optimum of the complete master LP --- obtained
without ever enumerating $\mathcal{R}$.

\end{document}
